\documentclass[12pt,letterpaper]{article}

% ── Paquetes ──────────────────────────────────────────────────────────────────
\usepackage[utf8]{inputenc}
\usepackage[T1]{fontenc}
\usepackage[spanish]{babel}
\usepackage{amsmath,amssymb,amsthm}
\usepackage{geometry}
\usepackage{hyperref}
\usepackage{listings}
\usepackage{xcolor}
\usepackage{booktabs}
\usepackage{array}
\usepackage{enumitem}
\usepackage{titlesec}
\usepackage{fancyhdr}
\usepackage{lmodern}
\usepackage{microtype}
\usepackage{parskip}

\geometry{margin=2.5cm}

% ── Colores y estilo de código ────────────────────────────────────────────────
\definecolor{codebg}{RGB}{245,245,245}
\definecolor{codeframe}{RGB}{200,200,200}
\definecolor{kwcolor}{RGB}{0,0,180}
\definecolor{strcolor}{RGB}{163,21,21}
\definecolor{cmtcolor}{RGB}{0,128,0}
\definecolor{numcolor}{RGB}{9,134,88}

\lstdefinestyle{python}{
  language=Python,
  backgroundcolor=\color{codebg},
  frame=single,
  rulecolor=\color{codeframe},
  basicstyle=\ttfamily\small,
  keywordstyle=\color{kwcolor}\bfseries,
  stringstyle=\color{strcolor},
  commentstyle=\color{cmtcolor}\itshape,
  numberstyle=\tiny\color{numcolor},
  numbers=left,
  numbersep=8pt,
  breaklines=true,
  breakatwhitespace=true,
  showstringspaces=false,
  tabsize=4,
  captionpos=b,
  xleftmargin=1.5em,
  framexleftmargin=1.5em,
}

% ── Encabezado y pie ──────────────────────────────────────────────────────────
\pagestyle{fancy}
\fancyhf{}
\rhead{\small C\'alculo Estoc\'astico para Finanzas}
\lhead{\small Universidad Panamericana}
\cfoot{\thepage}

% ── Teoremas ──────────────────────────────────────────────────────────────────
\newtheorem{theorem}{Teorema}[section]
\newtheorem{proposition}[theorem]{Proposici\'on}
\newtheorem{definition}{Definici\'on}[section]
\newtheorem{remark}{Observaci\'on}[section]

% ── Título ────────────────────────────────────────────────────────────────────
\title{
  \LARGE\bfseries C\'alculo Estoc\'astico para Finanzas\\[0.5em]
  \large Compendio Completo: Laboratorios, Tarea, Proyecto y Ejercicios
}
\author{
  Mauricio Olivares\\[0.3em]
  \small Profesor: Dr.\ Julio C\'esar Galindo\\
  \small Universidad Panamericana
}
\date{Verano 2026}

% ══════════════════════════════════════════════════════════════════════════════
\begin{document}
% ══════════════════════════════════════════════════════════════════════════════

\maketitle
\thispagestyle{empty}

\tableofcontents
\newpage

% ══════════════════════════════════════════════════════════════════════════════
\part{Laboratorios}
% ══════════════════════════════════════════════════════════════════════════════

\section{Laboratorio 1: Estimaci\'on de un Movimiento Browniano Geom\'etrico}

\subsection{Marco te\'orico}

El \textbf{Movimiento Browniano Geom\'etrico (GBM)} es el modelo est\'andar para la
evoluci\'on del precio de un activo financiero. La ecuaci\'on diferencial estoc\'astica es:

\begin{equation}
  dS_t = \mu S_t \, dt + \sigma S_t \, dB_t
\end{equation}

cuya soluci\'on expl\'icita, obtenida mediante la F\'ormula de It\^o, es:

\begin{equation}
  S_t = S_0 \exp\!\left[\left(\mu - \frac{\sigma^2}{2}\right)t + \sigma B_t\right]
\end{equation}

Los par\'ametros se calibran a partir de log-retornos diarios
$r_t = \ln(S_t / S_{t-1})$:

\begin{align}
  \hat{\mu}    &= 252 \cdot \bar{r} + \tfrac{1}{2}(252 \cdot s^2) \\
  \hat{\sigma} &= \sqrt{252} \cdot s
\end{align}

donde $s$ es la desviaci\'on est\'andar muestral de los log-retornos diarios.

\subsection{Objetivos}
\begin{enumerate}
  \item Descargar precios ajustados de AAPL (2020--2023).
  \item Estimar $\mu_b$ y $\sigma_b$ con log-retornos hist\'oricos.
  \item Simular 100 trayectorias GBM y comparar con el precio hist\'orico real.
  \item Graficar bandas de percentiles (5\textdegree, 50\textdegree, 95\textdegree).
\end{enumerate}

\subsection{C\'odigo}

\begin{lstlisting}[style=python, caption={Laboratorio 1 -- GBM con AAPL}]
import pandas as pd
import numpy as np
import matplotlib.pyplot as plt

try:
    import yfinance as yf
except ImportError:
    import os; os.system('pip install yfinance'); import yfinance as yf

ticker_symbol = 'AAPL'
df = yf.download(ticker_symbol, start='2020-01-01',
                 end='2023-12-31', progress=False)
if isinstance(df.columns, pd.MultiIndex):
    df.columns = df.columns.droplevel(1)
price_col = 'Adj Close' if 'Adj Close' in df.columns else 'Close'
df[price_col] = pd.to_numeric(df[price_col], errors='coerce')
df = df.dropna(subset=[price_col])

df['Log_Return'] = np.log(df[price_col]).diff()
df = df.dropna(subset=['Log_Return'])

mb = df['Log_Return'].mean()
sb = df['Log_Return'].std()

N       = 252
mu_b    = N * mb + 0.5 * (N * sb**2)
sigma_b = np.sqrt(N) * sb
print(f"Mu_b    = {mu_b:.6f}")
print(f"Sigma_b = {sigma_b:.6f}")

S0              = df[price_col].iloc[-1]
num_steps       = len(df)
num_simulations = 100
dt              = 1.0
daily_drift     = (mu_b - 0.5 * sigma_b**2) / N
daily_diffusion = sigma_b / np.sqrt(N)

simulated_paths = np.zeros((num_steps, num_simulations))
simulated_paths[0] = S0
np.random.seed(42)
for t in range(1, num_steps):
    Z = np.random.normal(0, 1, num_simulations)
    simulated_paths[t] = (simulated_paths[t-1]
        * np.exp(daily_drift*dt + daily_diffusion*np.sqrt(dt)*Z))

lower = np.percentile(simulated_paths, 5,  axis=1)
med   = np.percentile(simulated_paths, 50, axis=1)
upper = np.percentile(simulated_paths, 95, axis=1)

fig, ax = plt.subplots(figsize=(14, 6))
ax.plot(df.index, df[price_col], color='blue', lw=2,
        label='Historico real (AAPL)')
for i in range(min(20, num_simulations)):
    ax.plot(df.index, simulated_paths[:, i], color='gray', alpha=0.15)
ax.plot(df.index, med,   color='red',   ls='--', label='Percentil 50')
ax.plot(df.index, lower, color='green', ls=':',  label='Bandas 5-95')
ax.plot(df.index, upper, color='green', ls=':')
ax.set_title('AAPL vs. Simulaciones GBM', fontsize=14)
ax.set_xlabel('Fecha'); ax.set_ylabel('Precio (USD)')
ax.legend(); ax.grid(True, ls='--', alpha=0.6)
plt.tight_layout(); plt.show()
\end{lstlisting}

\subsection{Conclusiones}

El GBM ofrece una aproximaci\'on razonable de la tendencia y rango general de
movimiento de precios, pero presenta limitaciones importantes:

\begin{enumerate}
  \item \textbf{Volatilidad no constante (\textit{volatility clustering}):} los
    mercados exhiben periodos de calma seguidos de estallidos de alta volatilidad,
    incompatibles con $\sigma$ constante.
  \item \textbf{Colas pesadas:} los log-retornos reales tienen colas m\'as gruesas
    que una normal; el GBM subestima la probabilidad de eventos extremos.
  \item \textbf{Saltos discontinuos:} noticias y eventos geopo\'iticos producen
    saltos que el GBM no puede modelar.
  \item \textbf{Deriva constante:} el retorno esperado real cambia con el tiempo
    seg\'un fundamentos econ\'omicos.
\end{enumerate}

Modelos m\'as avanzados: volatilidad estoc\'astica (Heston), saltos-difusi\'on
(Merton), GARCH.

% ─────────────────────────────────────────────────────────────────────────────
\section{Laboratorio 2: Volatilidad Realizada y Ventanas M\'oviles}
% ─────────────────────────────────────────────────────────────────────────────

\subsection{Marco te\'orico}

La \textbf{volatilidad realizada anualizada} con ventana m\'ovil de $w$ d\'ias:

\begin{equation}
  \hat{\sigma}_t^{(w)} = \sqrt{252} \cdot
    \operatorname{std}\!\left\{r_{t-w+1}, \ldots, r_t\right\}
\end{equation}

\begin{table}[h!]
\centering
\begin{tabular}{@{}cll@{}}
\toprule
Ventana & Equivalencia & Interpretaci\'on \\
\midrule
21 d\'ias  & 1 mes burs\'atil    & Alta sensibilidad, mucho ruido \\
63 d\'ias  & 1 trimestre       & Balance sensibilidad/suavizado \\
126 d\'ias & 1 semestre        & Visi\'on estructural, efecto rezago \\
\bottomrule
\end{tabular}
\caption{Ventanas m\'oviles analizadas.}
\end{table}

\textbf{Volatility Clustering} (Mandelbrot; Engle, Premio Nobel -- modelos ARCH):
grandes cambios en precios son seguidos por grandes cambios, y peque\~nos por peque\~nos.

\subsection{C\'odigo}

\begin{lstlisting}[style=python, caption={Laboratorio 2 -- Volatilidad realizada}]
import pandas as pd
import numpy as np
import matplotlib.pyplot as plt
import seaborn as sns

try:
    import yfinance as yf
except ImportError:
    import os; os.system('pip install yfinance'); import yfinance as yf

sns.set_theme(style='whitegrid')

ticker = 'SPY'
datos  = yf.download(ticker, start='2014-01-01',
                     end='2024-01-01', progress=False)
if isinstance(datos.columns, pd.MultiIndex):
    datos.columns = datos.columns.droplevel(1)

precios = datos['Close'].copy()
log_ret = np.log(precios / precios.shift(1)).dropna()

ventanas = [21, 63, 126]
df_vol   = pd.DataFrame(index=log_ret.index)
for v in ventanas:
    df_vol[f'Vol_{v}d'] = log_ret.rolling(window=v).std() * np.sqrt(252)

fig, ax = plt.subplots(figsize=(15, 7))
ax.plot(df_vol['Vol_21d']  * 100, label='21d (corto)',
        color='#e74c3c', alpha=0.6, lw=1.2)
ax.plot(df_vol['Vol_63d']  * 100, label='63d (medio)',
        color='#3498db', alpha=0.8, lw=1.8)
ax.plot(df_vol['Vol_126d'] * 100, label='126d (largo)',
        color='#2c3e50', alpha=0.9, lw=2.5)
ax.axvspan(pd.to_datetime('2020-03-01'), pd.to_datetime('2020-06-01'),
           color='red', alpha=0.12)
ax.axvspan(pd.to_datetime('2022-01-01'), pd.to_datetime('2022-12-31'),
           color='orange', alpha=0.10)
ax.set_title(f'Volatilidad Realizada Anualizada -- {ticker}', fontsize=14)
ax.set_xlabel('Fecha'); ax.set_ylabel('Volatilidad (%)')
ax.legend(); ax.grid(True, ls='--', alpha=0.6)
plt.tight_layout(); plt.show()
\end{lstlisting}

\subsection{Conclusiones}

\begin{enumerate}
  \item \textbf{Reg\'imenes de mercado:} COVID-19 (mar--may 2020) llev\'o la ventana
    de 21 d\'ias por encima del 60\,\%; el ciclo de alzas de la Fed (2022) estableci\'o
    una meseta sostenida del 25--30\,\%.
  \item \textbf{Volatility clustering validado:} los picos forman bloques
    temporales extendidos, no eventos aislados.
  \item El supuesto de $\sigma$ constante del GBM b\'asico queda invalidado
    emp\'iricamente.
\end{enumerate}

% ─────────────────────────────────────────────────────────────────────────────
\section{Laboratorio 3: Valuaci\'on Monte Carlo de una Opci\'on Europea}
% ─────────────────────────────────────────────────────────────────────────────

\subsection{Marco te\'orico}

Bajo la medida neutral al riesgo $\mathbb{Q}$:

\begin{equation}
  C_0 = e^{-rT}\,\mathbb{E}^\mathbb{Q}[(S_T - K)^+]
\end{equation}

El m\'etodo Monte Carlo simula $N$ precios terminales:

\begin{equation}
  S_T^{(i)} = S_0 \exp\!\left[\left(r - \frac{\sigma^2}{2}\right)T
    + \sigma\sqrt{T}\,Z^{(i)}\right], \quad Z^{(i)} \sim \mathcal{N}(0,1)
\end{equation}

El estimador y su error est\'andar:

\begin{equation}
  \hat{C}_0 = \frac{e^{-rT}}{N}\sum_{i=1}^N (S_T^{(i)} - K)^+,
  \qquad
  \mathrm{SE} = \frac{e^{-rT}\,s}{\sqrt{N}}
\end{equation}

El benchmark anal\'itico es \textbf{Black-Scholes}:

\begin{align}
  C_0^{\mathrm{BS}} &= S_0\,\mathcal{N}(d_1) - Ke^{-rT}\,\mathcal{N}(d_2) \\
  d_1 &= \frac{\ln(S_0/K) + (r + \sigma^2/2)T}{\sigma\sqrt{T}}, \quad
  d_2 = d_1 - \sigma\sqrt{T}
\end{align}

\subsection{C\'odigo}

\begin{lstlisting}[style=python, caption={Laboratorio 3 -- Monte Carlo europeo}]
import numpy as np
import pandas as pd
from scipy.stats import norm
import datetime

try:
    import yfinance as yf
except ImportError:
    import os; os.system('pip install yfinance'); import yfinance as yf

ticker = 'SPY'
start  = (datetime.date.today()-datetime.timedelta(days=365*2)).strftime('%Y-%m-%d')
end    = datetime.date.today().strftime('%Y-%m-%d')
data   = yf.download(ticker, start=start, end=end, progress=False)
if isinstance(data.columns, pd.MultiIndex):
    data.columns = data.columns.droplevel(1)
S0      = data['Close'].iloc[-1].item()
lr      = np.log(data['Close']/data['Close'].shift(1)).dropna()
sigma_b = lr.std().item() * np.sqrt(252)

K               = S0 * 1.05
T               = 1.0
r               = 0.03
num_simulations = 100000

Z       = np.random.standard_normal(num_simulations)
ST      = S0 * np.exp((r - 0.5*sigma_b**2)*T + sigma_b*np.sqrt(T)*Z)
payoffs = np.maximum(ST - K, 0)

mc_price = np.exp(-r*T) * np.mean(payoffs)
se       = np.std(payoffs)/np.sqrt(num_simulations) * np.exp(-r*T)
ci_low   = mc_price - 1.96*se
ci_high  = mc_price + 1.96*se

d1 = (np.log(S0/K)+(r+0.5*sigma_b**2)*T)/(sigma_b*np.sqrt(T))
d2 = d1 - sigma_b*np.sqrt(T)
bs_price = S0*norm.cdf(d1) - K*np.exp(-r*T)*norm.cdf(d2)

print(f"MC Price : {mc_price:.4f}  SE: {se:.6f}")
print(f"IC 95%   : [{ci_low:.4f}, {ci_high:.4f}]")
print(f"BS Price : {bs_price:.4f}  Diff: {mc_price-bs_price:.6f}")
\end{lstlisting}

\subsection{Conclusiones}

\begin{itemize}
  \item El precio MC converge al precio Black-Scholes; la diferencia tiende a
    cero con $N \to \infty$ (convergencia del TCL).
  \item El error est\'andar decrece en proporci\'on a $1/\sqrt{N}$.
  \item Monte Carlo es valioso para opciones ex\'oticas sin soluci\'on anal\'itica
    cerrada.
\end{itemize}

% ─────────────────────────────────────────────────────────────────────────────
\section{Laboratorio 4: Opci\'on Americana con Longstaff--Schwartz}
% ─────────────────────────────────────────────────────────────────────────────

\subsection{Marco te\'orico}

Una opci\'on americana satisface el problema de \textbf{paro \'optimo}:

\begin{equation}
  V^{Am}(t, s) = \sup_{\tau \in \mathcal{T}_{t,T}}
    \mathbb{E}^\mathbb{Q}\!\left[e^{-r(\tau-t)} g(S_\tau) \,\Big|\, S_t = s\right]
\end{equation}

Como la europea corresponde a $\tau = T$, se tiene $V^{Am} \geq V^{Eu}$.
La desigualdad puede ser estricta para puts: ejercer anticipadamente recibe $K$
reinvertible a tasa $r$.

\textbf{Algoritmo de Longstaff--Schwartz (LSM):} regresi\'on hacia atr\'as desde $T$.
En cada paso se estima el valor de continuaci\'on mediante regresi\'on polinomial
sobre las trayectorias in-the-money:

\begin{equation}
  \hat{C}(S_t) = \sum_{k} \alpha_k L_k(S_t)
\end{equation}

Se ejerce cuando $g(S_t) > \hat{C}(S_t)$.

\subsection{C\'odigo}

\begin{lstlisting}[style=python, caption={Laboratorio 4 -- Longstaff-Schwartz}]
import numpy as np
import pandas as pd
from scipy.stats import norm

try:
    import yfinance as yf
except ImportError:
    import os; os.system('pip install yfinance'); import yfinance as yf

TICKER = 'SPY'
data   = yf.download(TICKER, start='2023-01-01',
                     end='2024-01-01', progress=False)
if isinstance(data.columns, pd.MultiIndex):
    data.columns = data.columns.droplevel(1)
S0      = data['Close'].iloc[-1].item()
lr      = np.log(data['Close']/data['Close'].shift(1))
sigma_b = lr.std().item() * np.sqrt(252)

K, T, r   = S0*0.95, 1.0, 0.05
num_paths = 10000
num_steps = 100
dt        = T / num_steps

np.random.seed(42)
paths    = np.zeros((num_steps+1, num_paths))
paths[0] = S0
for t in range(1, num_steps+1):
    Z        = np.random.normal(0, 1, num_paths)
    paths[t] = paths[t-1] * np.exp(
        (r - 0.5*sigma_b**2)*dt + sigma_b*np.sqrt(dt)*Z)

cash_flow     = np.maximum(K - paths[-1], 0)
exercise_time = np.full(num_paths, num_steps)

for t in range(num_steps-1, 0, -1):
    itm   = np.where(K - paths[t] > 0)[0]
    if len(itm) == 0: continue
    S_itm = paths[t][itm]
    y     = cash_flow[itm] * np.exp(-r*dt*(exercise_time[itm]-t))
    coef  = np.polyfit(S_itm, y, 2)
    cont  = np.polyval(coef, S_itm)
    iv    = np.maximum(K - S_itm, 0)
    ex    = itm[iv > cont]
    cash_flow[ex]     = iv[iv > cont]
    exercise_time[ex] = t

american_price = np.mean(cash_flow * np.exp(-r*dt*exercise_time))

d1 = (np.log(S0/K)+(r+0.5*sigma_b**2)*T)/(sigma_b*np.sqrt(T))
d2 = d1 - sigma_b*np.sqrt(T)
european_price = K*np.exp(-r*T)*norm.cdf(-d2) - S0*norm.cdf(-d1)

print(f"Put Americana (LSM) : {american_price:.4f}")
print(f"Put Europea   (BS)  : {european_price:.4f}")
print(f"Prima Americanidad  : {american_price-european_price:.4f}")
\end{lstlisting}

\subsection{Conclusiones}

\begin{itemize}
  \item La \textbf{prima de americanidad} cuantifica el valor de la flexibilidad
    de ejercicio temprano.
  \item Para calls sin dividendos la prima es cero (\'optimo no ejercer antes).
  \item LSM converge con $N \geq 10{,}000$ trayectorias y grado polinomial
    $\geq 2$; es el est\'andar en la industria.
\end{itemize}

% ─────────────────────────────────────────────────────────────────────────────
\section{Laboratorio 5: Opci\'on Parisina y Reloj de Barrera}
% ─────────────────────────────────────────────────────────────────────────────

\subsection{Marco te\'orico}

Una \textbf{opci\'on parisina down-and-out} se desactiva si el precio permanece
\emph{consecutivamente} por debajo de la barrera $H$ durante $D$ d\'ias.

El espacio de estados requiere el reloj de barrera $C_t$:

\begin{equation}
  V = V(t,\, S_t,\, C_t)
\end{equation}

La EDP incluye el t\'ermino $\partial V/\partial C$ cuando $S_t < H$:

\begin{equation}
  V_t + rS\,V_S + \tfrac{1}{2}\sigma^2 S^2 V_{SS}
    + \mathbf{1}_{\{S < H\}}\,V_C - rV = 0
\end{equation}

\textbf{Sensibilidades:}
\begin{itemize}
  \item $D \uparrow$ $\Rightarrow$ menor prob.\ de knock-out $\Rightarrow$
    precio sube (converge al europeo puro cuando $D \to \infty$).
  \item $H \uparrow$ $\Rightarrow$ mayor prob.\ de knock-out $\Rightarrow$
    precio baja.
\end{itemize}

\subsection{C\'odigo}

\begin{lstlisting}[style=python, caption={Laboratorio 5 -- Opcion parisina}]
import numpy as np
import pandas as pd
import matplotlib.pyplot as plt

try:
    import yfinance as yf
except ImportError:
    import os; os.system('pip install yfinance'); import yfinance as yf

ticker  = 'SPY'
data    = yf.download(ticker, start='2020-01-01',
                      end='2023-01-01', progress=False)
if isinstance(data.columns, pd.MultiIndex):
    data.columns = data.columns.droplevel(1)
S0      = data['Close'].iloc[-1].item()
lr      = np.log(data['Close']/data['Close'].shift(1)).dropna()
sigma_b = lr.std().item() * np.sqrt(252)

K, T, r         = S0*0.95, 1.0, 0.02
num_simulations = 100000
num_steps       = 252
dt              = T / num_steps
H               = S0 * 0.90
D_values        = [5, 10, 20, 30]

np.random.seed(42)
paths       = np.zeros((num_simulations, num_steps+1))
paths[:, 0] = S0
Z_all       = np.random.normal(0, 1, (num_simulations, num_steps))
for t in range(num_steps):
    paths[:, t+1] = paths[:, t] * np.exp(
        (r - 0.5*sigma_b**2)*dt + sigma_b*np.sqrt(dt)*Z_all[:, t])

def parisian_price(K, T, r, H, D_val, paths):
    n_sim   = paths.shape[0]
    payoffs = np.zeros(n_sim)
    for i in range(n_sim):
        path   = paths[i]
        ko     = False
        consec = 0
        for t in range(1, len(path)):
            if path[t] < H:
                consec += 1
                if consec >= D_val: ko = True; break
            else:
                consec = 0
        if not ko:
            payoffs[i] = max(0, K - path[-1])
    price = np.mean(payoffs) * np.exp(-r*T)
    se    = np.std(payoffs)/np.sqrt(n_sim) * np.exp(-r*T)
    return price, se

prices, errors = [], []
for D in D_values:
    p, e = parisian_price(K, T, r, H, D, paths)
    prices.append(p); errors.append(e)
    print(f"D={D:2d} dias -> Precio: {p:.4f}  SE: {e:.4f}")

fig, ax = plt.subplots(figsize=(9, 5))
ax.errorbar(D_values, prices, yerr=errors, fmt='-o', capsize=5,
            color='steelblue', ecolor='lightcoral', markerfacecolor='navy')
ax.set_title('Precio Parisino vs. Duracion D del Reloj', fontsize=13)
ax.set_xlabel('D (dias consecutivos bajo H)')
ax.set_ylabel('Precio de la Opcion')
ax.grid(True, ls='--', alpha=0.7)
plt.tight_layout(); plt.show()
\end{lstlisting}

\subsection{Conclusiones}

\begin{itemize}
  \item La \emph{memoria temporal} distingue a la opci\'on parisina de la barrera
    est\'andar.
  \item El espacio de estados debe ampliarse a $(t, S_t, C_t)$; sin $C_t$, dos
    estados financieramente distintos ser\'ian indistinguibles.
  \item Las opciones parisinas modelan escenarios donde la activaci\'on requiere
    confirmaci\'on sostenida (p.ej., regulaciones con umbral mantenido varios d\'ias).
\end{itemize}

% ══════════════════════════════════════════════════════════════════════════════
\part{Tarea: Procesos de Poisson, Cadenas de Markov y Martingalas Discretas}
% ══════════════════════════════════════════════════════════════════════════════

\section{Procesos de Poisson}

\begin{definition}
Un \textbf{proceso de Poisson} $\{N(t)\}_{t \geq 0}$ con tasa $\lambda > 0$
satisface:
\begin{enumerate}
  \item $N(0) = 0$.
  \item Incrementos independientes y estacionarios.
  \item $P(N(t) = k) = \dfrac{(\lambda t)^k e^{-\lambda t}}{k!}$,\quad
    $E[N(t)] = \lambda t$,\quad $\mathrm{Var}[N(t)] = \lambda t$.
\end{enumerate}
\end{definition}

\begin{proposition}[Ley condicional]
  Para $0 < s < t$:\quad $N(s)\,|\,N(t)=n \;\sim\; \mathrm{Binomial}(n,\, s/t)$.
\end{proposition}

\subsection*{Ejercicio 1.1 --- Llegadas de \'ordenes en una mesa de FX
  ($\lambda = 10$ \'ord/hora)}

\begin{enumerate}[label=(\alph*)]
  \item $P(N(0.5)=0) = e^{-10 \times 0.5} = e^{-5} \approx 0.0067$.

    \textit{Implicaci\'on:} periodos de iliquidez donde el riesgo de inventario
    aumenta.

  \item Por independencia de incrementos en intervalos disjuntos:
    \[
      P(N_{10:30\text{-}11}=3)\cdot P(N_{11:30\text{-}12}=7)
      = \frac{e^{-5}5^3}{3!}\cdot\frac{e^{-5}5^7}{7!} \approx 0.0146
    \]

  \item $E[N(8)] = 80$,\quad $\mathrm{SD} = \sqrt{80} \approx 8.94$.

    No es prudente dimensionar solo con la media por la alta volatilidad del
    flujo.

  \item $\mathrm{SD}(N(T)) = \sqrt{\lambda T} \propto \sqrt{T}$, relacionado
    con la \textbf{regla de la ra\'iz del tiempo} del VaR:
    $\mathrm{VaR}_{10d} \approx \sqrt{10}\,\mathrm{VaR}_{1d}$.
\end{enumerate}

\subsection*{Ejercicio 1.4 --- Ley condicional}

\[
  N(s)\,|\,N(t)=n \;\sim\; \mathrm{Binomial}\!\left(n,\,\frac{s}{t}\right)
\]
\[
  E[N(s)\,|\,N(t)=n] = n\frac{s}{t}, \qquad
  \mathrm{Var}[N(s)\,|\,N(t)=n] = n\frac{s}{t}\left(1-\frac{s}{t}\right)
\]

\subsection*{Ejercicio 1.8 --- Poisson Compuesto (Cram\'er-Lundberg)}

$E[X(t)] = \lambda t\,E[Y]$,\quad $\mathrm{Var}[X(t)] = \lambda t\,E[Y^2]$.

Con $\lambda=3$ siniestros/mes y $E[Y]=10{,}000$ MXN:
\[
  E[X(12)] = 3 \times 12 \times 10{,}000 = 360{,}000 \text{ MXN}
\]

\section{Cadenas de Markov}

\begin{definition}
  Una \textbf{cadena de Markov} satisface
  $P(X_{n+1}=j\,|\,X_n=i_n,\ldots,X_0=i_0) = p_{ij}$.
  La \textbf{distribuci\'on estacionaria} $\pi$ cumple $\pi P = \pi$,
  $\sum_i \pi_i = 1$.
\end{definition}

\subsection*{Ejercicio 2.1 --- Clasificaci\'on y distribuci\'on estacionaria}

Para una matriz de transici\'on $P$ con estados $\{1,2,3\}$ (ratings crediticios),
se resuelve $\pi P = \pi$ para obtener la distribuci\'on de largo plazo.

\subsection*{Ejercicio 2.4 --- Absorci\'on y riesgo crediticio}

Se resuelve $u_i = \sum_j p_{ij} u_j$ para calcular la probabilidad de absorci\'on
en el estado de default ($R_2$) o solvencia ($R_1$).

\subsection*{Ejercicio 2.6 --- Rating con Default (estado absorbente)}

$P^3(1,4)$ da la probabilidad de default a 3 a\~nos para un emisor con rating
inicial AAA.

\section{Martingalas Discretas}

\begin{definition}
  Una sucesi\'on $\{M_n\}$ es \textbf{martingala} respecto a $\{\mathcal{F}_n\}$ si
  $M_n$ es $\mathcal{F}_n$-medible, $E[|M_n|]<\infty$ y
  $E[M_{n+1}\,|\,\mathcal{F}_n] = M_n$ c.s.
\end{definition}

\subsection*{Ejercicio 3.1 --- Modelo CRR}

Bajo la medida de riesgo neutro:
\[
  q = \frac{1+r-d}{u-d}
\]
El precio descontado $\tilde{S}_n = S_n/(1+r)^n$ es martingala bajo
$\mathbb{Q}$, garantizando la ausencia de arbitraje.

\subsection*{Ejercicio 3.3 --- Desigualdad de Doob}

\[
  P\!\left(M^* \geq \lambda\right) \leq \frac{E[M_N]}{\lambda}
\]

usando $\tau = \min\{k : M_k \geq \lambda\}$.

\subsection*{Ejercicio 3.5 --- Criterio de Kelly}

Se maximiza $G(f) = E[\log(W_n/W_0)]$. La soluci\'on $f^*$ optimiza el crecimiento
de largo plazo. Sobreinvertir ($f > f_{\max}$) lleva a $W_n \to 0$ c.s.

\subsection*{Ejercicio 3.8 --- FTAP y Valuaci\'on}

\begin{itemize}
  \item \textbf{Primer Teorema Fundamental:} no hay arbitraje $\iff$ existe una
    medida de martingala equivalente (EMM).
  \item \textbf{Segundo Teorema:} completitud del mercado $\iff$ unicidad de la
    EMM.
\end{itemize}

% ══════════════════════════════════════════════════════════════════════════════
\part{Proyecto Final: Factores y Estilos de Inversi\'on --- AAPL 2019--2024}
% ══════════════════════════════════════════════════════════════════════════════

\section{Introducci\'on y Pregunta Central}

\begin{quote}
  \textit{\'Puede describirse el comportamiento de AAPL mediante exposiciones a
  estilos o factores emp\'iricos?}
\end{quote}

\section{Marco Te\'orico}

\subsection{CAPM}

\begin{equation}
  E(R_i) = R_f + \beta_i\,[E(R_m) - R_f]
\end{equation}

$\alpha$ captura el retorno no explicado; en equilibrio deber\'ia ser cero.

\subsection{Modelo Fama-French 3 Factores (FF3)}

\begin{equation}
  R_i - R_f = \alpha + b_1(\text{Mkt-RF}) + b_2\,\text{SMB}
    + b_3\,\text{HML} + \varepsilon_i
\end{equation}

\begin{itemize}
  \item \textbf{SMB} (Small Minus Big): prima de tama\~no.
  \item \textbf{HML} (High Minus Low): prima de valor. Carga negativa
    $\Rightarrow$ perfil \textit{growth}.
\end{itemize}

\section{Datos}

\begin{itemize}
  \item Activo: AAPL (NASDAQ), febrero 2019 -- diciembre 2024,
    \textbf{71 observaciones mensuales}.
  \item Factores FF3 y $R_f$: Kenneth French Data Library.
  \item Estimaci\'on: OLS con errores robustos HC3 (White, 1980).
\end{itemize}

\section{Estad\'isticas Descriptivas}

\begin{table}[h!]
\centering
\begin{tabular}{@{}lrrrrr@{}}
\toprule
Variable & Media (\%) & Desv.\ Est. & M\'in.\ (\%) & M\'ax.\ (\%) & Sharpe \\
\midrule
AAPL   &  2.210 & 8.398 & $-14.71$ & 21.35 & 0.263 \\
Mkt-RF &  1.464 & 5.230 & $-13.37$ & 13.66 & 0.280 \\
SMB    & $-0.291$ & 2.721 & $-7.10$ & 7.47 & --- \\
HML    &  0.075 & 3.836 & $-9.89$  & 9.81 & --- \\
RF     &  0.202 & 0.175 & 0.01    & 0.47 & --- \\
\bottomrule
\end{tabular}
\caption{Estad\'isticas descriptivas -- retornos mensuales (\%).}
\end{table}

\section{Resultados}

\subsection{CAPM}

\begin{equation}
  \widehat{\mathrm{AAPL}_{exc}} = 0.0003 + 1.4192\cdot(\text{Mkt-RF}) + \varepsilon
  \qquad R^2 = 76.4\,\%
\end{equation}

$\beta = 1.4192$ ($p < 0.01$): por cada 1\,\% de retorno del mercado, AAPL
rinde $\approx 1.42$\,\% adicional. $\alpha$ no significativo ($p = 0.953$).

\subsection{Fama-French 3 Factores}

\begin{equation}
  \widehat{\mathrm{AAPL}_{exc}} = -0.0004 + 1.4594\,(\text{Mkt-RF})
    - 0.0022\,\text{SMB} - 0.0029\,\text{HML} + \varepsilon
  \qquad R^2 = 78.2\,\%
\end{equation}

\subsection{Comparaci\'on de Modelos}

\begin{table}[h!]
\centering
\begin{tabular}{@{}lrrr@{}}
\toprule
M\'etrica & CAPM & FF3 & $\Delta$ \\
\midrule
$R^2$ & 76.38\,\% & 78.15\,\% & +1.77 pp \\
AIC   & $-450.40$ & $-457.10$ & FF3 mejor \\
BIC   & $-445.88$ & $-448.05$ & FF3 mejor \\
Test $F$ incremental & --- & $F=4.155$ & $p=0.0198$ \\
\bottomrule
\end{tabular}
\caption{CAPM vs.\ Fama-French 3F.}
\end{table}

\section{Conclusiones del Proyecto}

\begin{enumerate}
  \item El comportamiento de AAPL se describe satisfactoriamente mediante
    exposiciones a factores sistem\'aticos (FF3 explica el 78.2\,\% de la
    varianza mensual).
  \item Factor de mercado dominante: $b_1 \approx 1.46$, $p<0.01$.
  \item Perfil factorial: \textbf{large-cap} ($b_2<0$), \textbf{growth}
    ($b_3<0$), \textbf{alta beta} --- coherente con mega-capitalizaci\'on
    tecnol\'ogica (\$3B USD).
  \item FF3 mejora al CAPM en $R^2$, AIC y BIC; el test $F$ incremental
    confirma la mejora ($p=0.020$).
  \item $\alpha$ no significativo en ning\'un modelo $\Rightarrow$ consistente
    con mercados eficientes.
\end{enumerate}

\section{C\'odigo del Proyecto}

\begin{lstlisting}[style=python, caption={Proyecto Final -- CAPM y FF3 para AAPL}]
import numpy as np
import pandas as pd
import matplotlib.pyplot as plt

try:
    import yfinance as yf
except ImportError:
    import os; os.system('pip install yfinance'); import yfinance as yf

try:
    from statsmodels.regression.linear_model import OLS
    from statsmodels.tools import add_constant
    import statsmodels.api as sm
except ImportError:
    import os; os.system('pip install statsmodels')
    from statsmodels.regression.linear_model import OLS
    from statsmodels.tools import add_constant
    import statsmodels.api as sm

try:
    import pandas_datareader.data as web
except ImportError:
    import os; os.system('pip install pandas-datareader')
    import pandas_datareader.data as web

aapl = yf.download('AAPL', start='2019-02-01',
                   end='2024-12-31', progress=False)
if isinstance(aapl.columns, pd.MultiIndex):
    aapl.columns = aapl.columns.droplevel(1)
aapl_monthly = aapl['Close'].resample('ME').last()
aapl_ret     = aapl_monthly.pct_change().dropna()
aapl_ret.index = aapl_ret.index.to_period('M')

ff3 = web.DataReader('F-F_Research_Data_Factors',
                     'famafrench', start='2019-01', end='2024-12')[0] / 100
ff3.index = ff3.index.to_period('M')

df = pd.concat([aapl_ret.rename('AAPL'), ff3], axis=1).dropna()
df['AAPL_exc'] = df['AAPL'] - df['RF']

X_capm   = add_constant(df['Mkt-RF'])
res_capm = OLS(df['AAPL_exc'], X_capm).fit(cov_type='HC3')
print("=== CAPM ===")
print(res_capm.summary2())

X_ff3   = add_constant(df[['Mkt-RF','SMB','HML']])
res_ff3 = OLS(df['AAPL_exc'], X_ff3).fit(cov_type='HC3')
print("=== Fama-French 3F ===")
print(res_ff3.summary2())

print(f"R2 CAPM: {res_capm.rsquared:.4f} | R2 FF3: {res_ff3.rsquared:.4f}")
print(f"AIC CAPM: {res_capm.aic:.2f}    | AIC FF3: {res_ff3.aic:.2f}")
\end{lstlisting}

% ══════════════════════════════════════════════════════════════════════════════
\part{20 Ejercicios Resueltos: Browniano e Integral de It\^o}
% ══════════════════════════════════════════════════════════════════════════════

\section*{Ejercicio 1 --- Incrementos de Browniano}
\addcontentsline{toc}{section}{Ejercicio 1 -- Incrementos de Browniano}

Para $0 \leq s < t \leq u < v$, usando $K(a,b)=\min\{a,b\}$:

\[
  \mathrm{Cov}(B_t-B_s,\, B_v-B_u) = K(t,v)-K(t,u)-K(s,v)+K(s,u)
  = t - t - s + s = 0
\]

Al ser gaussianos con covarianza nula $\Rightarrow$ independientes. $\blacksquare$

\section*{Ejercicio 2 --- Ganancia de una estrategia discreta}
\addcontentsline{toc}{section}{Ejercicio 2 -- Ganancia de una estrategia discreta}

$H_0=1,\; H_1=2,\; H_2=-1$:

\[
  (H\cdot B)_T = 1\cdot(B_{t_1}-B_0) + 2\cdot(B_{t_2}-B_{t_1})
    + (-1)\cdot(B_T-B_{t_2}) = -B_{t_1} + 3B_{t_2} - B_T
\]

Ganancia acumulada; $H_i$ se fija antes del siguiente incremento
(no anticipaci\'on). $\blacksquare$

\section*{Ejercicio 3 --- Isometr\'ia de It\^o}
\addcontentsline{toc}{section}{Ejercicio 3 -- Isometria de Ito}

$H_t = a\,\mathbf{1}_{(0,T/2]} + b\,\mathbf{1}_{(T/2,T]}$:

\[
  E\!\left[\!\left(\int_0^T H_t\,dB_t\right)^{\!2}\right]
  = a^2\cdot\frac{T}{2} + b^2\cdot\frac{T}{2}
  = E\!\left[\int_0^T H_t^2\,dt\right] \qquad \checkmark \qquad \blacksquare
\]

\section*{Ejercicio 4 --- Variaci\'on cuadr\'atica}
\addcontentsline{toc}{section}{Ejercicio 4 -- Variacion cuadratica}

Con partici\'on uniforme de $[0,T]$ con malla $T/n$:

\[
  E[V_n] = T, \quad
  \mathrm{Var}[V_n] = \frac{2T^2}{n} \to 0
  \quad \Rightarrow \quad V_n \xrightarrow{P} T = [B]_T \qquad \blacksquare
\]

\section*{Ejercicio 5 --- F\'ormula de It\^o para el logaritmo}
\addcontentsline{toc}{section}{Ejercicio 5 -- Formula de Ito para el logaritmo}

$f(x) = \log x$:

\[
  d(\log S_t) = \frac{dS_t}{S_t} + \frac{1}{2}\!\left(-\frac{1}{S_t^2}\right)
    (\sigma S_t)^2\,dt
  = \left(\mu - \frac{\sigma^2}{2}\right)dt + \sigma\,dB_t \qquad \blacksquare
\]

\section*{Ejercicio 6 --- Distribuci\'on bajo GBM}
\addcontentsline{toc}{section}{Ejercicio 6 -- Distribucion bajo GBM}

\[
  \log\frac{S_t}{S_0} = \left(\mu-\frac{\sigma^2}{2}\right)t + \sigma B_t
  \;\sim\; \mathcal{N}\!\!\left(\left(\mu-\frac{\sigma^2}{2}\right)t,\;\sigma^2 t\right)
\]

$S_t$ sigue distribuci\'on log-normal. El t\'ermino $-\sigma^2/2$ es la
\textbf{correcci\'on de It\^o}. $\blacksquare$

\section*{Ejercicio 7 --- Volatilidad realizada}
\addcontentsline{toc}{section}{Ejercicio 7 -- Volatilidad realizada}

\[
  r_t = \log(S_t/S_{t-1}), \quad
  s = \sqrt{\frac{1}{n-1}\sum(r_t-\bar{r})^2}, \quad
  \sigma_{\text{anual}} = \sqrt{252}\cdot s \qquad \blacksquare
\]

\section*{Ejercicio 8 --- Valuaci\'on de una opci\'on europea}
\addcontentsline{toc}{section}{Ejercicio 8 -- Valuacion de una opcion europea}

\[
  C_0 = e^{-rT}\,E^\mathbb{Q}[(S_T-K)^+]
\]

La delta $\Delta_t = \partial V/\partial S$ es la cobertura. $\blacksquare$

\section*{Ejercicio 9 --- Put americana vs.\ put europea}
\addcontentsline{toc}{section}{Ejercicio 9 -- Put americana vs put europea}

\[
  V^{Am}(t,s) = \sup_{\tau\in\mathcal{T}_{t,T}}
    E^\mathbb{Q}[e^{-r(\tau-t)}g(S_\tau)\,|\,S_t=s]
  \;\geq\; V^{Eu}(t,s) \qquad \blacksquare
\]

\section*{Ejercicio 10 --- Opci\'on parisina: estado ampliado}
\addcontentsline{toc}{section}{Ejercicio 10 -- Opcion parisina: estado ampliado}

El precio requiere $V = V(t, S_t, C_t)$. Sin $C_t$, dos estados
financieramente distintos ser\'ian indistinguibles. $\blacksquare$

\section*{Ejercicio 11 --- Martingala $M_t = B_t^2 - t$}
\addcontentsline{toc}{section}{Ejercicio 11 -- Martingala Mt = Bt2 - t}

\[
  E[M_t\,|\,\mathcal{F}_s] = B_s^2 + (t-s) - t = B_s^2 - s = M_s
  \quad\checkmark \qquad \blacksquare
\]

\section*{Ejercicio 12 --- Martingala exponencial}
\addcontentsline{toc}{section}{Ejercicio 12 -- Martingala exponencial}

$Z_t = e^{\theta B_t - \theta^2 t/2}$:

\[
  E[Z_t\,|\,\mathcal{F}_s]
  = Z_s \cdot e^{\theta^2(t-s)/2} \cdot e^{-\theta^2(t-s)/2}
  = Z_s \quad\checkmark \qquad \blacksquare
\]

$Z_T = dQ/dP$ es la densidad de Radon-Nikodym (Teorema de Girsanov).

\section*{Ejercicio 13 --- Precio de mercado del riesgo}
\addcontentsline{toc}{section}{Ejercicio 13 -- Precio de mercado del riesgo}

\[
  \lambda = \frac{\mu - r}{\sigma}
\]

Con $B_t^\mathbb{Q} = B_t^\mathbb{P} + \lambda t$, el drift pasa de $\mu$ a
$r$ bajo $\mathbb{Q}$. $\blacksquare$

\section*{Ejercicio 14 --- Precio descontado como martingala}
\addcontentsline{toc}{section}{Ejercicio 14 -- Precio descontado como martingala}

\[
  d(e^{-rt}S_t) = \sigma\,S_t^*\,dB_t^\mathbb{Q}
\]

Solo t\'ermino de difusi\'on $\Rightarrow$ drift nulo $\Rightarrow$
\textbf{martingala bajo $\mathbb{Q}$}. $\blacksquare$

\section*{Ejercicio 15 --- Paridad put--call}
\addcontentsline{toc}{section}{Ejercicio 15 -- Paridad put-call}

\[
  \boxed{C_0 - P_0 = S_0 - Ke^{-rT}} \qquad \blacksquare
\]

\section*{Ejercicio 16 --- Delta y Gamma}
\addcontentsline{toc}{section}{Ejercicio 16 -- Delta y Gamma}

\[
  \Delta V \approx \Delta\cdot\Delta S + \tfrac{1}{2}\Gamma(\Delta S)^2
    + \Theta\,\Delta t
\]

$\Delta = \partial V/\partial S$: cobertura lineal.
$\Gamma = \partial^2 V/\partial S^2$: curvatura; $\Gamma>0$ beneficia de la
volatilidad. $\blacksquare$

\section*{Ejercicio 17 --- PDE de Black--Scholes}
\addcontentsline{toc}{section}{Ejercicio 17 -- PDE de Black-Scholes}

\begin{equation}
  \boxed{V_t + \frac{1}{2}\sigma^2 S^2 V_{SS} + rS\,V_S - rV = 0,
    \quad V(T,S) = h(S)} \qquad \blacksquare
\end{equation}

\section*{Ejercicio 18 --- Frontera de ejercicio en put americana}
\addcontentsline{toc}{section}{Ejercicio 18 -- Frontera de ejercicio en put americana}

Condiciones de \textit{smooth pasting} en $S^*(t)$:

\[
  V(t, S^*(t)) = K - S^*(t), \qquad
  \frac{\partial V}{\partial S}(t, S^*(t)) = -1
\]

$S^*(t)$ decreciente; $S^*(T) = K$. $\blacksquare$

\section*{Ejercicio 19 --- Sensibilidad de opci\'on parisina down-and-out}
\addcontentsline{toc}{section}{Ejercicio 19 -- Sensibilidad de opcion parisina}

\begin{itemize}
  \item $H\uparrow$: mayor prob.\ de knock-out $\Rightarrow$ \textbf{precio baja}.
  \item $D\uparrow$: menor prob.\ de desactivaci\'on $\Rightarrow$ \textbf{precio sube}.
\end{itemize}
Se invierte para knock-in. $\blacksquare$

\section*{Ejercicio 20 --- Volatilidad hist\'orica vs.\ impl\'icita}
\addcontentsline{toc}{section}{Ejercicio 20 -- Volatilidad historica vs implicita}

\begin{table}[h!]
\centering
\begin{tabular}{@{}lll@{}}
\toprule
 & Volatilidad hist\'orica & Volatilidad impl\'icita \\
\midrule
Direcci\'on  & Backward-looking & Forward-looking \\
Estimaci\'on & $\sqrt{252}\cdot\mathrm{std}\{r_t\}$
             & Se despeja de $C_{\text{mercado}} = BS(\sigma)$ \\
\bottomrule
\end{tabular}
\caption{Volatilidad hist\'orica vs.\ impl\'icita.}
\end{table}

Su diferencia es la \textbf{variance risk premium}. $\blacksquare$

% ══════════════════════════════════════════════════════════════════════════════
\part*{Bibliograf\'ia}
\addcontentsline{toc}{part}{Bibliografia}
% ══════════════════════════════════════════════════════════════════════════════

\begin{enumerate}
  \item Fama, E.F.\ \& French, K.R.\ (1993). Common risk factors in the returns
    on stocks and bonds. \textit{Journal of Financial Economics}, 33(1), 3--56.
  \item Longstaff, F.A.\ \& Schwartz, E.S.\ (2001). Valuing American options by
    simulation. \textit{Review of Financial Studies}, 14(1), 113--147.
  \item Shreve, S.E.\ (2004). \textit{Stochastic Calculus for Finance II}.
    Springer.
  \item Sharpe, W.F.\ (1964). Capital asset prices. \textit{Journal of Finance},
    19(3), 425--442.
  \item White, H.\ (1980). A heteroskedasticity-consistent covariance matrix
    estimator. \textit{Econometrica}, 48(4), 817--838.
\end{enumerate}

\end{document}
